%%%%%%%%%%%%%%%%%%%%%%%%%%%%%%%
\chapter{Redesign Compressible Flow Solver}
\label{chap:redesign-cfs}

This chapter collects the parts of \code{CompressibleFlowSolverRedesign} that
concern the \code{SolverCore} framework rather than its discretisation. Its
equation systems, \code{EulerCFE} and \code{NavierStokesCFE}, derive from
\code{SolverCore::UnsteadySystem}.

%%%%%%%%%%%%%%%%%%%%%%%%%%%%%%%
\section{Filters}
\label{sec:redesign-cfs-filters}

The solver keeps the default field registry of \code{EquationSystem}
(Chapter~\ref{chap:solvercore}): \code{m\_fields}, the conservative variables
density, momentum and total energy, are the solution role and the one
physical group. Checkpoint, AverageFields, HistoryPoints and Mean therefore
work unchanged, and their files contain the conservative variables.

The compressible solution is conservative, so component slices of
\code{m\_fields} cannot be used as velocity or pressure views. Supporting the
fluid diagnostics requires the equation system to provide backend-native
primitive workspaces and populate

\begin{itemize}
\item velocity as momentum divided by density;
\item pressure through the configured equation of state; and
\item density for density-weighted diagnostics.
\end{itemize}

These values should be computed lazily once per filter phase and shared by all
consumers. A natural extension is a lazy physical-role provider in
\code{FieldRegistry}, with the resolved view cached for the lifetime of that
registry. The conversion kernel and its workspace belong to the compressible
solver (or to a compressible primitive-state service), not to the filters.
This preserves the backend-independent filter layer and avoids a host
conversion of conservative fields.

Fluid diagnostics need more than the primitive roles. A kinetic-energy
diagnostic of variable-density flow integrates density-weighted quantities,
so it requires the density role as well as the velocity. A boundary-force
diagnostic requires a constitutive contract as well as primitive pressure:
the Navier--Stokes equations may use a dynamic, possibly
temperature-dependent, viscosity and a dilatational stress term, and Euler
flow has no viscous traction. The equation system should therefore publish a
viscous-stress (or boundary traction) view or service computed by its
diffusion model. Treating momentum gradients as velocity gradients would be
incorrect.

Tests of such diagnostics need analytic variable-density fields, and
inviscid and viscous cases using the same equation-of-state and transport
models as the solver. A negative test should confirm that an Euler run
rejects a requested viscous contribution.
